% Part: quantified-modal-logic
% Chapter: semantics
% Section: comparisons

\documentclass[../../../include/open-logic-section]{subfiles}

\begin{document}

\olfileid{qml}{sem}{com}

\olsection{Comparison Betweeen Classical vs Free Quantified Modal Logics}

\begin{center}
Comparing validity across classical quantified modal logic (\Log{CQML})
and free quantified modal logics: negative (\Log{NFQML}) and 
positive (\Log{PFQML}).
\smallskip
\begin{tabular}[t]{p{.20\textwidth}p{.30\textwidth}ccc}
    \hline
    !!^{sentence} & English trans. & \multicolumn{3}{c}{Valid in:} \\
        &       & \Log{CQML} & \Log{NFQML} & \Log{PFQML} \\
    \hline 
    \multicolumn{2}{l}{Necessity of Existence} \\
    $\Box\lforall[x][\lfrexists x]$ 
        & Necessarily, everything exists 
        & $\surd$ & $\surd$ & $\surd$ \\
    $\lforall[x][\Box\lfrexists x]$ 
        & Everything exists necessarily
        & $\surd$ & $\times$ & $\times$ \\
    $\lfrexists \Obj{a}$ 
        & $\Obj{a}$ exists necessarily
        & $\surd$ & $\times$ & $\times$ \\
    \multicolumn{2}{l}{Conditional Necessity of Identity} \\
    $\eq[\Obj{a}][\Obj{b}]\lif\Box(\lfrexists\Obj{a}\lif \eq[\Obj{a}][\Obj{b}])$
        & If $\Obj{a}$ is identical to $\Obj{b}$ then necessarily if $\Obj{a}$ exists it is identical to $\Obj{a}$
        & $\surd$ & $\surd$ & $\surd$ \\
    $\lforall[x][\lforall[y][(\eq[x][y]\lif\Box\eq[x][y])]]$
        & If something is identical to a thing then necessarily it is identical to that thing
        & $\surd$ & $\times$ & $\surd$ \\
    \multicolumn{2}{l}{Unconditional Necessity of Identity} \\
        $\eq[\Obj{a}][\Obj{b}]\lif\Box\eq[\Obj{a}][\Obj{b}]$
        & If $\Obj{a}$ is identical to $\Obj{b}$ then necessarily $\Obj{a}$ is identical to $\Obj{b}$
        & $\surd$ & $\times$ & $\surd$ \\
    $\lforall[x][\lforall[y][(\eq[x][y]\lif\Box\eq[x][y])]]$
        & If something is identical to a thing then necessarily it is identical to that thing
        & $\surd$ & $\times$ & $\surd$ \\
    \multicolumn{2}{l}{Identity} \\
    $\Box\lforall[x][\eq[x][x]]$
        & Necessarily, everything is itself
        & $\surd$ & $\surd$ & $\surd$ \\
    $\Box\eq[\Obj{a}][\Obj{a}]$
        & $\Obj{a}$ is $\Obj{a}$
        & $\surd$ & $\times$ & $\surd$ \\
    \multicolumn{2}{l}{Atomic predication} \\
    $\Box(\Atom{\Obj{P}}{\Obj{a}}\lif\lfrexists\Obj{a})$ 
        & Necessarily, if $\Obj{a}$ is $\Obj{P}$, then $\Obj{a}$ exists 
        & $\surd$ & $\surd$ & $\times$ \\
    \multicolumn{2}{l}{Negated atomic predication} \\
    $\Box(\lnot \Atom{\Obj{P}}{\Obj{a}}\lif\lfrexists\Obj{a})$ 
        & Necessarily, if it's not the case that $\Obj{a}$ is $\Obj{P}$, then $\Obj{a}$ exists 
        & $\surd$ & $\times$ & $\times$ \\
    \hline
\end{tabular}
\smallskip
\end{center}

\emph{Note on the existence predicate $\lfrexists$}. In free logics
it is equivalent to $\lexists[v][\eq[\ldots][v]]$. In first-order
logic with identity we haven't introduced it, but it would be
equivalent to that to. Thus were first-order logic is concerned,
simply treat $\lfrexists c$ as an abbreviation of
 $\lexists[v][\leq[c][v]]$.
Thus:
\begin{itemize}
    \item $\lfrexists \Obj{a}$ is 
    equivalent to $\lexists[x][\eq[\Obj{a}][x]]$
    \item $\lforall[x][\lfrexists x]$ is 
    equivalent $\lforall[x][\lexists[y][\eq[x][y]]]$
    \item $\lfrexists\Obj{a}\lif\Subst{!A}{\Obj{a}}{x}$ is 
    equivalent $\lexists[y][\eq[\Obj{a}][y]]\lif\Subst{!A}{\Obj{a}}{x}$
\end{itemize}

\end{document}
